\clearpage
\section{Perencanaan, Deskripsi dan Hasil Uji Perangkat Lunak}

\subsection{Rencana Pengujian}

Pengujian fungsional dilakukan dengan pendekatan \textit{black-box} terhadap seluruh skenario \textit{use case} yang diimplementasikan (\texttt{UC-01} s.d.\ \texttt{UC-13}). Pengujian turut mencakup verifikasi persistensi tab pasca-aksi pada dasbor Penjual serta Admin, audit protokol transfer hiperteks (\textit{Hypertext Transfer Protocol}) untuk metode permintaan \texttt{GET} dan \texttt{POST}, serta pengujian unit/fitur bawaan Laravel. Eksekusi pengujian otomatis dijalankan pada tanggal 19 November 2026 memakai \textit{Playwright} pada peramban \textit{Chromium} terhadap klon basis data terisolasi (\texttt{jajanria\_test}) agar data usaha mikro, kecil, dan menengah (UMKM) pada basis data utama tetap aman. Daftar kasus uji dirangkum pada Tabel~\ref{tab:rencana-pengujian}.

\renewcommand{\arraystretch}{1.22}
\begin{longtable}{|C{1.55cm}|L{\dimexpr\textwidth-1.55cm-2.0cm-1.7cm-2.1cm-2.3cm-12\tabcolsep-7\arrayrulewidth\relax}|C{2.0cm}|C{1.7cm}|C{2.1cm}|L{2.3cm}|}
\caption{Rencana Pengujian Fungsional Sistem JajanRia}
\label{tab:rencana-pengujian} \\
\hline
\rowcolor{TableHeadBg}
\textbf{Nomor Uji} & \textbf{Nama Kasus Uji} & \textbf{Kebutuhan} & \textbf{Jenis Uji} & \textbf{Tanggal Uji} & \textbf{Nama Penguji} \\
\hline
\endfirsthead
\hline
\rowcolor{TableHeadBg}
\textbf{Nomor Uji} & \textbf{Nama Kasus Uji} & \textbf{Kebutuhan} & \textbf{Jenis Uji} & \textbf{Tanggal Uji} & \textbf{Nama Penguji} \\
\hline
\endhead
\hline
\endfoot
KU-001 & Autentikasi Admin dan Penjual valid serta penolakan akun nonaktif atau sandi salah & F001 & \textit{Black-box} & 19/11/2026 & Tim Penguji TRPL-104 \\
\hline
KU-002 & Registrasi mandiri mitra Penjual baru melalui formulir pendaftaran mitra & F001 & \textit{Black-box} & 19/11/2026 & Tim Penguji TRPL-104 \\
\hline
KU-003 & Pencarian kata kunci, penyaring katalog, dan halaman detail lapak UMKM & F002--F005 & \textit{Black-box} & 19/11/2026 & Tim Penguji TRPL-104 \\
\hline
KU-004 & \textit{Checkout} keranjang, kalkulasi harga di server, dan penolakan menu habis & F006 & \textit{Black-box} & 19/11/2026 & Tim Penguji TRPL-104 \\
\hline
KU-005 & Unggah foto bukti pembayaran transfer saat \textit{checkout} dan via \textit{History Pesanan} & F007 & \textit{Black-box} & 19/11/2026 & Tim Penguji TRPL-104 \\
\hline
KU-006 & Verifikasi bukti pembayaran terpusat oleh Admin (Setujui \textit{Lunas} \& Tolak via modal alasan) & F008 & \textit{Black-box} & 19/11/2026 & Tim Penguji TRPL-104 \\
\hline
KU-007 & Penjual menyelesaikan pesanan \textit{Self Pick-up} dan larangan sebelum status \textit{lunas} & F009 & \textit{Black-box} & 19/11/2026 & Tim Penguji TRPL-104 \\
\hline
KU-008 & Pemicu tombol ``Beri Penilaian'' (\textit{selesai} dan \textit{lunas}), pengiriman ulasan, serta balasan Penjual & F010--F012 & \textit{Black-box} & 19/11/2026 & Tim Penguji TRPL-104 \\
\hline
KU-009 & Pengelolaan menu, profil lapak UMKM, dan status ketersediaan menu Penjual & F013--F014 & \textit{Black-box} & 19/11/2026 & Tim Penguji TRPL-104 \\
\hline
KU-010 & Verifikasi calon Penjual, proteksi hapus kategori berproduk aktif, dan moderasi ulasan Admin & F015--F017, F019 & \textit{Black-box} & 19/11/2026 & Tim Penguji TRPL-104 \\
\hline
KU-011 & Komunikasi layanan \textit{Customer Support} dua arah antara Pengunjung dan Admin Support & F018 & \textit{Black-box} & 19/11/2026 & Tim Penguji TRPL-104 \\
\hline
KU-012 & Otorisasi akses Panel Administrator, Dasbor Penjual, dan isolasi kepemilikan ulasan & F020 & \textit{Black-box} & 19/11/2026 & Tim Penguji TRPL-104 \\
\hline
\end{longtable}

\subsection{Deskripsi dan Hasil Uji}

% Macro kotak hasil uji standar Halaman 41 Format laporan Polibatam
\newcommand{\printtestcase}[9]{%
\needspace{9cm}
\noindent
\renewcommand{\arraystretch}{1.20}
\begin{tabularx}{\textwidth}{|L{3.2cm}|Y|Y|}
\hline
\rowcolor{TableHeadBg}
\textbf{Identifikasi} & \multicolumn{2}{l|}{\textbf{#1}} \\
\hline
\textbf{Nama Kasus Uji} & \multicolumn{2}{>{\raggedright\arraybackslash}p{\dimexpr\textwidth-3.2cm-4\tabcolsep-3\arrayrulewidth\relax}|}{#2} \\
\hline
\textbf{Deskripsi} & \multicolumn{2}{>{\raggedright\arraybackslash}p{\dimexpr\textwidth-3.2cm-4\tabcolsep-3\arrayrulewidth\relax}|}{#3} \\
\hline
\textbf{Kondisi Awal} & \multicolumn{2}{>{\raggedright\arraybackslash}p{\dimexpr\textwidth-3.2cm-4\tabcolsep-3\arrayrulewidth\relax}|}{#4} \\
\hline
\textbf{Tanggal Pengujian} & \multicolumn{2}{l|}{#5} \\
\hline
\textbf{Penguji} & \multicolumn{2}{l|}{#6} \\
\hline
\rowcolor{RowShade}
\multicolumn{3}{|c|}{\textbf{Skenario Pengujian}} \\
\hline
\multicolumn{3}{|>{\raggedright\arraybackslash}p{\dimexpr\textwidth-2\tabcolsep-2\arrayrulewidth\relax}|}{#7} \\
\hline
\rowcolor{TableHeadBg}
\textbf{Hasil yang Diharapkan} & \textbf{Hasil Pengamatan} & \textbf{Kesimpulan} \\
\hline
#8 & #9 & \textbf{\color{JajanGreen}Sesuai (PASS)} \\
\hline
\end{tabularx}
\vspace{0.35cm}
}

Rincian skenario, hasil yang diharapkan, hasil pengamatan aktual, dan kesimpulan untuk seluruh 12 kasus uji fungsional disajikan pada tabel identifikasi di bawah ini.

\printtestcase{KU-001}{Autentikasi Admin dan Penjual Valid serta Penolakan Sandi Salah / Akun Nonaktif}
{Menguji proses masuk memakai surel dan kata sandi valid maupun penolakan kata sandi literal yang bukan milik akun, serta pemilih akun cepat aktif.}
{Akun Admin, Penjual aktif, dan Pengunjung tersedia di basis data uji.}
{19 November 2026}{Tim Penguji TRPL-104 (Playwright \& Audit HTTP)}
{\begin{uclist}
  \item Kirim permintaan masuk memakai surel aktif dengan kata sandi tidak valid (\texttt{wrongpass} dan literal \texttt{'password'}).
  \item Klik tombol akses cepat peran Penjual pada halaman masuk.
  \item Masuk menggunakan kredensial Administrator valid.
\end{uclist}}
{Percobaan sandi salah ditolak ke halaman masuk; tombol cepat Penjual memilih lapak aktif; Admin masuk ke Panel Administrator.}
{Percobaan sandi literal \texttt{'password'} dialihkan kembali ke halaman masuk (\texttt{HTTP 302}); akses cepat Penjual masuk ke lapak aktif Mie Aceh Bang Din; Admin masuk ke Panel Administrator.}

\printtestcase{KU-002}{Registrasi Mandiri Mitra Penjual Baru melalui Formulir Pendaftaran}
{Menguji pengiriman formulir pendaftaran mitra UMKM dengan seluruh bidang selaras beserta validasinya.}
{Pengguna belum masuk akun dan membuka halaman pendaftaran mitra.}
{19 November 2026}{Tim Penguji TRPL-104 (Playwright E2E)}
{\begin{uclist}
  \item Isi nama pemilik, nama usaha (\texttt{Bakso Budi Kampus}), surel baru, nomor kontak, area (\texttt{Tower A}), alamat lokasi, dan kata sandi.
  \item Tekan tombol daftar mitra (\texttt{DAFTAR MITRA}).
  \item Periksa halaman tujuan dan status verifikasi pengguna pada Panel Administrator.
\end{uclist}}
{Pengguna dialihkan ke halaman masuk dengan spanduk sukses, serta calon mitra tercatat pada antrean verifikasi Admin.}
{Pengalihan ke halaman masuk berhasil dengan spanduk ``Pendaftaran Mitra UMKM berhasil!'', dan data \texttt{Bakso Budi Kampus} muncul di tab verifikasi Admin.}

\printtestcase{KU-003}{Pencarian Kata Kunci, Penyaring Katalog, dan Halaman Detail Lapak}
{Menguji tampilan katalog beranda, penyaringan, serta pemunculan lapak terverifikasi baru di halaman publik.}
{Mitra baru (\texttt{Bakso Budi Kampus}) telah disetujui Admin pada basis data uji.}
{19 November 2026}{Tim Penguji TRPL-104 (Playwright E2E)}
{\begin{uclist}
  \item Buka halaman beranda katalog sebagai Pengunjung dan cari lapak terverifikasi baru.
  \item Buka halaman detail lapak dan periksa kelompok menu serta status ketersediaan produk.
\end{uclist}}
{Lapak terverifikasi baru tampil di katalog publik beranda dan rincian menu tampil pada halaman detail lapak.}
{Lapak \texttt{Bakso Budi Kampus} tampil di katalog utama; halaman detail lapak memuat daftar menu per kategori dan penanda menu habis.}

\printtestcase{KU-004}{Checkout Keranjang, Kalkulasi Harga Server, dan Blokir Menu Habis}
{Menguji konsistensi nominal pada modal \textit{checkout}, validasi harga asli server, penyimpanan catatan pembeli, dan penolakan menu berstatus habis.}
{Pengunjung masuk akun dan membuka halaman detail lapak dengan salah satu menu tersedia dan satu menu berstatus habis.}
{19 November 2026}{Tim Penguji TRPL-104 (Playwright \& Audit HTTP)}
{\begin{uclist}
  \item Tambahkan 1 porsi menu (Rp22.000) ke keranjang lalu buka modal \textit{checkout} dan periksa total tagihan pada ringkasan modal.
  \item Kirim permintaan \textit{checkout} dengan nominal harga manipulatif (Rp1) untuk 2 porsi menu Rp22.000.
  \item Kirim permintaan \textit{checkout} untuk menu berstatus habis.
\end{uclist}}
{Modal menampilkan \texttt{Rp 22.000}; server mengabaikan harga manipulatif dan menghitung total Rp44.000; menu habis ditolak (\texttt{HTTP 422}).}
{Terbaca \texttt{'Rp 22.000'} di modal; permintaan \textit{checkout} mengembalikan \texttt{HTTP 200} dengan total tercatat \texttt{44000}; menu habis ditolak \texttt{HTTP 422}.}

\printtestcase{KU-005}{Unggah Foto Bukti Pembayaran Saat Checkout dan melalui History Pesanan}
{Menguji dua skenario pembayaran: melampirkan bukti langsung saat \textit{checkout}, serta \textit{checkout} tanpa bukti lalu mengunggah dari \textit{History Pesanan}.}
{Pengunjung telah masuk akun pada sesi aktif.}
{19 November 2026}{Tim Penguji TRPL-104 (Playwright E2E)}
{\begin{uclist}
  \item Lakukan \textit{checkout} pada lapak pertama disertai berkas gambar bukti bayar (\texttt{WF-01}).
  \item Lakukan \textit{checkout} tanpa foto pada lapak kedua, lalu buka laci \textit{History Pesanan} dan unggah foto bukti transfer (\texttt{WF-02}).
\end{uclist}}
{Pesanan langsung berstatus pembayaran \textit{Menunggu Verifikasi} dan status pesanan \textit{Diproses}; unggah susulan dari History berhasil.}
{Pesanan tercatat di \textit{History Pesanan} dengan status \textit{Menunggu Verifikasi} dan \textit{Diproses}, serta unggah susulan berhasil.}

\printtestcase{KU-006}{Verifikasi Bukti Pembayaran Terpusat oleh Admin (Setujui dan Tolak via Modal)}
{Menguji pengubahan status pembayaran oleh Admin dari tab pembayaran, termasuk pembukaan modal alasan penolakan serta persistensi tab.}
{Terdapat pesanan berstatus menunggu verifikasi pada basis data uji.}
{19 November 2026}{Tim Penguji TRPL-104 (Playwright E2E)}
{\begin{uclist}
  \item Masuk sebagai Admin lalu buka tab Verifikasi Pembayaran pada Panel Administrator.
  \item Tekan tombol Tolak untuk membuka modal penolakan pembayaran, isi alasan penolakan, lalu simpan.
  \item Tekan tombol Setujui (\textit{Lunas}) pada pesanan uji berikutnya.
\end{uclist}}
{Modal penolakan terbuka normal, status pembayaran diperbarui menjadi \textit{ditolak}/\textit{lunas}, dan tampilan tetap bertahan pada tab Verifikasi Pembayaran.}
{Modal tolak terbuka mulus, kedua aksi berhasil mengubah status pembayaran dan posisi halaman tetap berada di tab Verifikasi Pembayaran.}

\printtestcase{KU-007}{Penjual Menyelesaikan Pesanan Self Pick-up setelah Pembayaran Lunas}
{Menguji penampilan rincian item dan catatan pembeli di tabel Penjual, penguncian tombol Selesai sebelum lunas, dan penyelesaian setelah lunas.}
{Terdapat pesanan belum lunas dan pesanan lunas milik lapak Mie Aceh Bang Din.}
{19 November 2026}{Tim Penguji TRPL-104 (Playwright \& Audit HTTP)}
{\begin{uclist}
  \item Buka tab Pesan \& Pesanan pada Dasbor Penjual dan periksa kolom rincian menu, catatan pembeli, serta tombol status.
  \item Lakukan upaya mengubah status pesanan menjadi \textit{selesai} pada pesanan yang belum \textit{lunas}.
  \item Klik tombol Selesai pada pesanan yang telah diverifikasi \textit{lunas} oleh Admin.
\end{uclist}}
{Nama menu dan catatan tampil; pesanan belum lunas terkunci dan ditolak server; pesanan lunas berubah \textit{selesai} dan tetap berada di tab Pesan \& Pesanan.}
{Nama menu serta catatan pembeli tampil; penolakan sebelum lunas tampil; pesanan lunas berubah menjadi \textit{selesai} dan tampilan tetap di tab Pesan \& Pesanan.}

\printtestcase{KU-008}{Pemicu Tombol ``Beri Penilaian'' (Selesai + Lunas), Kirim Ulasan, dan Balasan Penjual}
{Menguji pemicu penilaian dan penolakan ulasan belum lunas. Diuji pula penyimpanan ulasan di History serta blokir balasan ulasan tersembunyi (\texttt{UC-08}).}
{Terdapat pesanan belum lunas, pesanan \textit{selesai} dan \textit{lunas}, serta satu ulasan disembunyikan Admin.}
{19 November 2026}{Tim Penguji TRPL-104 (Playwright \& Audit HTTP)}
{\begin{uclist}
  \item Kirim permintaan pengiriman ulasan untuk pesanan yang pembayarannya belum \textit{lunas}.
  \item Buka \textit{History Pesanan} untuk pesanan \textit{selesai} dan \textit{lunas}, klik ``Beri Penilaian'', lalu kirim bintang 5 dan komentar.
  \item Masuk sebagai Penjual pada tab Rating \& Ulasan, balas ulasan tampil, lalu coba balas ulasan yang disembunyikan Admin (\texttt{UC-08}).
\end{uclist}}
{Permintaan menolak ulasan belum lunas (\texttt{HTTP 400}); tombol penilaian muncul saat \textit{selesai} serta \textit{lunas}; balasan ulasan tersembunyi ditolak.}
{Ulasan belum lunas ditolak \texttt{HTTP 400}; tombol ``Beri Penilaian'' muncul dan tetap di laci History; balasan pada ulasan disembunyikan ditolak pesan galat.}

\printtestcase{KU-009}{Pengelolaan Menu, Profil Lapak UMKM, dan Ketersediaan Menu Penjual}
{Menguji tambah/ubah/hapus/toggle produk dan pembaruan profil termasuk area luar sekitar kampus, serta persistensi tab Penjual.}
{Penjual aktif masuk ke Dasbor Penjual.}
{19 November 2026}{Tim Penguji TRPL-104 (Playwright E2E)}
{\begin{uclist}
  \item Pada tab Kelola Menu \& Katalog, tambah menu baru, klik ubah ketersediaan (\textit{Tersedia}/\textit{Habis}), sunting menu, lalu hapus menu.
  \item Pada tab Profil Lapak UMKM, perbarui profil lapak lalu tekan simpan.
\end{uclist}}
{Seluruh operasi dan simpan profil berhasil tersimpan di basis data serta mempertahankan posisi tab aksi masing-masing.}
{Seluruh 9 skenario persistensi tab Penjual (termasuk 4 aksi produk, profil, pesanan, dan ulasan + muat ulang F5) berstatus \textbf{PASS}.}

\printtestcase{KU-010}{Verifikasi Calon Penjual, Proteksi Hapus Kategori (\texttt{UC-10}), dan Moderasi Ulasan Admin}
{Menguji persetujuan calon Penjual baru, pencegahan penghapusan kategori berproduk aktif, serta moderasi ulasan oleh Admin.}
{Admin masuk ke Panel Administrator; terdapat calon mitra baru \texttt{Bakso Budi Kampus} dan kategori berproduk aktif.}
{19 November 2026}{Tim Penguji TRPL-104 (Playwright E2E)}
{\begin{uclist}
  \item Pada tab Verifikasi Seller Baru, tekan Setujui pada pengajuan calon Penjual baru.
  \item Pada tab Kelola Kategori Menu, tambah kategori uji baru, coba hapus kategori yang memiliki produk aktif (\texttt{UC-10}), lalu hapus kategori uji kosong.
  \item Pada tab Moderasi Ulasan, tekan Sembunyikan/Pulihkan pada ulasan.
\end{uclist}}
{Calon Penjual disetujui; hapus kategori berproduk aktif ditolak dengan peringatan; moderasi ulasan berhasil dan semua tetap di tab aksi.}
{Seluruh 10 skenario persistensi tab Admin (kategori, pembayaran, UMKM, moderasi, support, pengguna, verifikasi) berstatus \textbf{PASS}.}

\printtestcase{KU-011}{Komunikasi Layanan Customer Support Dua Arah}
{Menguji pengiriman tiket keluhan oleh Pengunjung, tampilan label \texttt{Admin Support} di log monitoring, serta balasan Admin kembali ke laci Pengunjung.}
{Pengunjung dan Admin dapat masuk pada konteks peramban masing-masing.}
{19 November 2026}{Tim Penguji TRPL-104 (Playwright E2E)}
{\begin{uclist}
  \item Pengunjung membuka laci \textit{Customer Support} lalu mengirim pesan pertanyaan bantuan.
  \item Admin membuka tab Customer Support pada Panel Administrator, membaca tiket masuk, lalu mengirim balasan via modal.
  \item Cek tampilan balasan pada laci Pengunjung dan label pengirim pada log \textit{Monitoring \& Statistik}.
\end{uclist}}
{Pesan tampil pada tab Support Admin, balasan Admin terbaca langsung di laci Pengunjung, dan log menampilkan label \texttt{Admin Support}.}
{Tiket terkirim, balasan Admin tampil di laci Pengunjung, dan label \texttt{'Admin Support'} tampil di monitoring.}

\printtestcase{KU-012}{Otorisasi Hak Akses Peran dan Isolasi Kepemilikan Ulasan}
{Menguji penolakan tamu maupun akses lintas peran ke Panel Administrator dan Dasbor Penjual, serta isolasi balasan ulasan antar-UMKM.}
{Pengguna berstatus tamu dan akun Pengunjung masuk tersedia untuk pengujian hak akses.}
{19 November 2026}{Tim Penguji TRPL-104 (Audit HTTP Backend)}
{\begin{uclist}
  \item Lakukan akses langsung ke Panel Administrator dan Dasbor Penjual tanpa sesi login.
  \item Lakukan akses ke Panel Administrator dan Dasbor Penjual menggunakan sesi berperan \texttt{pengunjung}.
  \item Kirim permintaan balasan ulasan sebagai Penjual lapak A terhadap ulasan milik lapak B.
\end{uclist}}
{Tamu dialihkan ke halaman masuk (\texttt{HTTP 302}); peran salah ditolak (\texttt{HTTP 403}); balasan ulasan lintas UMKM ditolak (\texttt{HTTP 404}).}
{Tamu dialihkan \texttt{HTTP 302} ke halaman masuk; Pengunjung mendapat \texttt{HTTP 403} pada Panel Admin dan Dasbor Penjual; balasan ulasan UMKM lain ditolak \texttt{HTTP 404}.}

\subsection{Rekapitulasi Pengujian Otomatis}

Selain 12 kasus uji fungsional di atas, pengujian unit dan fitur Laravel (\texttt{php artisan test}) dijalankan setelah seluruh perubahan kode selesai dan menghasilkan \textbf{2 lulus (\textit{2 passed, 2 assertions})} untuk \texttt{Tests\textbackslash{}Unit\textbackslash{}ExampleTest} dan \texttt{Tests\textbackslash{}Feature\textbackslash{}ExampleTest}. Pada seluruh pengujian peramban \textit{Playwright} atas alur kerja \texttt{WF-01..06}, pemeriksaan penyelarasan \texttt{EC-01..03}, serta 19 skenario persistensi tab (10 Admin + 9 Penjual), tercatat \textbf{0 galat konsol} dan \textbf{0 galat halaman}.
